\chapter{Implementació}
\label{ch:implementacio}

\section{Tecnologies utilitzades}
\subsection{Frontend}
% TODO: Next.js 16, React 19, TypeScript, Tailwind 4, Shadcn/UI. Justificació vs alternatives.

\subsection{Backend}
% TODO: Supabase (Postgres + pgvector + Auth), Edge Functions. Justificació vs Firebase/backend propi.

\subsection{Intel·ligència artificial}
% TODO: Claude Haiku 4.5 (chat), Gemini embedding-001 (768 dim). Per què Anthropic i no OpenAI/Gemini directe.

\subsection{Model Context Protocol}
% TODO: justificació del MCP com a estàndard 2026 (cites: Anthropic, AWS, Cloudflare, OWASP).

\section{Migració SQL per a MCP i agents}
% TODO: explicar la migració 20260419120000_mcp_tfg.sql (taules agent_events, tag_suggestions, índex HNSW).

\section{Servidor MCP}
% TODO: explicar src/app/api/mcp/route.ts. Decisions: Streamable HTTP stateless, Bearer-token PoC -> OAuth.

\section{Agents en segon pla}
% TODO: Edge Function per cada agent. Programació via pg_cron. Auditoria via agent_events.

\section{Decisions destacades}
% TODO: llista 5-10 decisions tècniques no òbvies amb raonament.

\section{Auditoria estructural del codi via graphify (2026-04-24)}
\label{sec:graphify-audit}

% Durant la sessió de documentació post-merge es va executar graphify
% (pipeline graph-RAG: AST + LLM-inference + Louvain clustering) sobre
% la totalitat del repositori per validar estructuralment el disseny
% descrit a \S\ref{ch:disseny}. El corpus: 132 fitxers, ~90.365 mots.
% Resultat: 380 nodes, 427 arestes, 77 comunitats detectades sense
% supervisió. Valor de la tècnica en el context del TFG: les troballes
% següents són verificables pel tribunal amb el `graphify-out/graph.html`
% interactiu inclòs al repositori, a diferència d'afirmacions
% arquitectòniques purament narratives.

\subsection{Separació Part~A / Part~B verificada estructuralment}
% Hipòtesi del disseny: plataforma base (notes, xat, auth) i extensió
% MCP són dos mòduls disjunts que comparteixen només el fet de viure
% al mateix projecte Next.js. La graph ho confirma: el node `GET()`
% de `src/app/api/mcp/route.ts` és l'única aresta cross-community de
% pes significatiu (betweenness = 0.022) entre el cluster de server
% actions de notes/chat i el cluster MCP. No hi ha altres edges
% estructurals que creuin la frontera. En termes pràctics: si s'hagués
% d'extreure el servidor MCP a un repositori separat, el tall seria
% net (un sol fitxer com a superfície d'API + import del
% `NotesService`).

\subsection{Consolidació d'auth-gates confirmada per degree centrality}
% `requireUser()` (src/actions/settings.ts, src/actions/notes.ts) i
% `requireChatAccess()` (src/actions/chats.ts) surten amb 8 arestes
% cadascuna. Totes les server actions que muten fan passada obligatòria
% per una d'aquestes dues helpers abans de tocar la BD. És un patró
% belt-and-braces sobre les polítiques RLS (ambdues capes han de
% coincidir perquè la mutation surti): primer l'auth de Supabase
% (cookies/JWT) via les helpers, després RLS al costat Postgres. El
% graph verifica que cap ruta de mutation s'ha escapat del patró.

\subsection{Re-descobriment automàtic dels grups arquitectònics}
% Graphify va generar 4 hyperedges de grup sense que se li donés cap
% instrucció sobre quins conceptes agrupar --- els va derivar per
% co-ocurrència i similitud semàntica dins el corpus:
%
% 1. "MCP security stack" --- MCP + Lethal Trifecta + OAuth 2.1 + RLS +
%    filtre de sortida (D3) + Promptfoo. Confidence EXTRACTED 0.90.
%    Confirma que la narrativa de \S\ref{ch:disseny} sobre seguretat
%    és coherent: els 6 conceptes apareixen encadenats a les mateixes
%    fonts, no són afegits retroactivament.
% 2. "Sis eines MCP" --- search\_notes, get\_note, create\_note,
%    update\_note, tag\_notes, summarise\_notes. Confidence EXTRACTED
%    1.00. Coincideix amb l'abast planificat al \S\ref{ch:requisits}.
% 3. "Fases QoL amb tema star/archive/delete" --- QoL-3 + QoL-7 +
%    commit e94aea5 + patró useOptimistic. Confidence INFERRED 0.85.
%    Captura un arc de polishing que no estava nominalment agrupat
%    però comparteix solució (soft-mutation + optimistic UI).
% 4. "Variants de paleta explorades" --- A graphite-electric,
%    B zinc-emerald, C navy-amber. Confidence EXTRACTED 1.00.
%
% El primer i el segon són la validació més forta que docs i codi
% parlen del mateix sistema: un algoritme no-supervisat ha
% reconstruït les estructures dissenyades.

\subsection{Mapa de comunitats 1:1 amb els mòduls planificats}
% De les 77 comunitats, 15 són "denses" (>4 nodes) i les 62 restants
% són single/double-node (fitxers utility, configs, shadcn primitives
% individuals). El long-tail no és soroll sinó evidència de
% baix acoplament: els components primitius viuen aïllats perquè
% realment no tenen dependències creuades.
%
% Les 15 comunitats grans mapegen directament als mòduls del disseny:
% server actions de notes/auth, lògica del sidebar del xat, backend
% MCP + agents, conceptes TFG Part~B + recerca, implementació del
% servidor MCP (JWT + NotesService), screenshots de baseline, paletes,
% settings view, NoteGrid handlers, tags manager, compose zone,
% shortcuts globals, i les famílies de primitives shadcn (Sheet,
% Select, Dialog). Cap comunitat creua dos mòduls del disseny.
% El clustering Louvain (modular, no supervisat) ha re-derivat la
% partició intencional.

\subsection{Fals positiu detectat i la seva rellevància metodològica}
% El node amb més centralitat bruta (degree 21) és `Select()`. És un
% error sistemàtic d'extracció: l'LLM infereix arestes `calls` cap a
% `Select()` des de qualsevol server action que conté la seqüència
% `.select()` (mètode del client de Supabase). Confon el component
% `<Select>` de shadcn/ui (src/components/ui/select.tsx) amb la
% cadena `supabase.from(...).select(...)` per coincidència léxica.
%
% L'episodi és útil per al capítol d'avaluació (vegeu
% \S\ref{sec:graphify-limitations}): totes aquestes arestes estan
% taggejades com a INFERRED, cosa que les fa auditables. El valor
% del pipeline híbrid AST+LLM no és produir un graph perfecte sinó
% produir-ne un amb errors *detectables* que una inspecció manual
% pot filtrar.

